\documentclass{article}
\usepackage[utf8]{inputenc}
\usepackage[T1]{fontenc}
\usepackage[margin=2cm]{geometry}
\usepackage{hyperref}
\sloppy
\begin{document}
\section*{Historical source}
Encyclopaedia Britannica scan, source row 220. Parse 3.5 extracted the text. Super 3.5 VL corrected it from an enlarged crop of each block. It can still contain recognition errors.
Page headers, footers, and page numbers are removed.

\subsection*{Untitled continuation (E0005)}
\textit{Continues an entry from an earlier page.}

FUR [512] FUR

and called by him Queen Charlotte's Sound, the latitude of which is 51 degrees north, longitude 128 west.

The snow Lark, Captain Peters, of 220 tons and 40 men, sailed from Macao in July 1786. Her destination was Kamtschatka (for which she was provided with a suitable cargo of arrack, tea, \&c.), Copper Islands, and the N. W. coast. Captain Peters was directed to make his passage between Japan and Corea, and examine the islands to the north of Japan, said to be inhabited by hairy people; which, if Captain Cook had lived, would not have been left to the French to determine. No account having been received of this vessel since her departure, there is every reason to fear she has perished.

In the beginning of 1786, two coppered vessels were fitted out at Bombay, under the direction of James Strange, Esq; who was himself a principal owner. These vessels were, the snow Captain Cook of 300 tons, and snow Experiment of 100 tons. They proceeded in company from the Malabar coast to Batavia; passed through the Straits of Macassar, where the Experiment was run upon a reef, and was obliged to haul ashore upon Borneo to repair; from thence they steered to the eastward of the Paloes Islands; made Sulphur Island; and arrived at Nootka the end of June following. From Nootka, where they left their surgeon's mate (Mackay) to learn the language and collect skins against their intended return (but who was brought away in the Imperial Eagle the following year), they proceeded along the coast to Queen Charlotte's Sound, of which they were the first discoverers; from thence in a direct course to Prince William's Sound. After some stay there, the Experiment proceeded to Macao (their vessels being provided with pafles by the governor-general of Goa); the Captain Cook endeavoured to get to Copper Island, but without success, being prevented by contrary winds.

Two coppered vessels were also fitted out by a society of gentlemen in Bengal, viz. the snow Nootka of 200 tons, and the snow Sea Otter of 100 tons, commanded by John Meares and William Tipping, lieutenants in the royal navy. The Nootka sailed in March 1786 from Bengal; came through the China Seas; touched at the Baffles, where they were very civilly treated by the Spaniards, who have taken possession of these islands; arrived at Oonalaska the beginning of August; found there a Russian galliot and some furriers; discovered accidentally near Cape Greville a new strait into Cook's River, 15 leagues wide and 30 long; saw some Russian hunters in a small bay between Cape Elizabeth and Cape Bear; and arrived in Prince William's Sound the end of September. They determined wintering in Snug Corner Cove, lat. 60. 30. in preference to going to the Sandwich Islands, which seem placed by Providence for the comfort and refreshment of the adventurers in this trade, and were frozen up in this gloomy and frightful spot from the end of November to the end of May. By the severity of the winter they lost their third and fourth mates, fourgon, boatwain, carpenter, and cooper, and twelve of the foremast men; and the remainder were so enfeebled as to be under the necessity of applying to the commanders of the King George and Queen Charlotte, who just at this time arrived in the sound, for some hands to assist in carrying the ve-

fell to the Sandwich Islands, where, giving over all further thoughts of trade, they determined (after getting a sea-fleet of fish off Cape Edgecumbe) immediately to proceed. The Nootka arrived at Macao the end of October 1787.

The Imperial Eagle, Captain Barkley, fitted out by a society of gentlemen at Olend, sailed from Olend the latter end of November 1786; went into the bay of All Saints; from thence, without touching any where, to the Sandwich Islands, and arrived at Nootka the beginning of June; from thence to the fourth, as far as 47° 30', in which space he discovered some good and spacious harbours. In the lat. of 47° 46', lost his second mate, purser, and two seamen, who were upon a trading party with the long-boat, and imprudently trusting themselves ashore unarmed, were cut off by the natives. This place seems to be the fame that Don Antonio Mourelle calls the Isla de los Dolores, where the Spaniards going ashore to water, were also attacked and cut off.

The King George of 320, and the Queen Charlotte of 200 tons, commanded by Captains Portlock and Dixon, who served under Captain Cook in his last voyage, were fitted out by a society of gentlemen in England, who obtained a privilege to trade to the north-west coast of America, from the South-Sea and East India companies.

Thofo vessels sailed from England the beginning of September 1785; touched at the Falkland Islands, Sandwich Islands, and arrived in Cook's River in the month of August. From thence, after collecting a few furs, they fleered in the end of September for Prince William's Sound, intending, it is said, to winter there; but were prevented entering, by heavy forms and extreme bad weather, which obliged them to bear away, and seek some other part of the coast to winter at. The forms and bad weather accompanied them till they arrived off Nootka Sound, when they were to raise the slows, that a cause some of us them; but though then now accomplishing their purpose, a Fifth items came on, and obliged them finally to bear away for the Sandwich Islands, where they remained the winter months; and returning again to the coast, arrived in Prince William's Sound the middle of May. The King George remained in Prince William's Sound; and during her stay, her long-boat discovered a new passage from the Sound into Cook's River. The Queen Charlotte proceeded along the coast to the fourth; looked into Behring's Bay, where the Ruffians have now a settlement; examined that part of the coast from 56 to 50, which was not seen by Captain Cook, and which consists of a cluster of islands, called by Captain Dixon Queen Charlotte's Islands, at a considerable distance from the Main, which is thus removed farther to the eastward than it was supposed to be; some part of the continent may, however, be seen from the east side of these islands; and it is probable, the distance does not any where exceed 50 leagues. On this elucidation, Hudson's Houle, lat. 53 long. 106 27 west, will not be more than 800 miles distant from that part of this coast in the same parallel. It is therefore not improbable, that the enterprising spirit of our Canadian furriers may penetrate to this coast (the communication with which is probably much facilitated by lakes or rivers), and add to the comforts and luxu


\section*{Source record}
\begin{description}
\item[Scan ID] \texttt{britannica:309f765b1ecc83e0a7f7762e44b4a927a7555443:train:220}
\item[Image file] \texttt{page-0220-f28b3445b235.jpg}
\item[Image SHA-256] \texttt{f28b3445b2354a93e2efbd2647cb2058637dea9d7c095bb9c2e239624ac95565}
\item[Dataset] NationalLibraryOfScotland/encyclopaedia\_britannica\_illustrated at revision \texttt{309f765b1ecc83e0a7f7762e44b4a927a7555443} (CC0-1.0)
\item[Extraction] model \texttt{nvidia/nemotron-3.5-parse}, raw reply \texttt{extract/raw/page-0220-f28b3445b235.json}
\item[Raw reply SHA-256] \texttt{499ec5efc188e0fd721d72914ae6c3f035bb2fc9dd9ddd5005b35874cca64e76}
\item[Correction] model \texttt{nvidia/Nemotron-3.5-Super-VL}, block crops, raw replies \texttt{correct/raw/crop\_nothink/page-0220-f28b3445b235-b*.json}; super\_vl: 9
\item[Running header] none
\end{description}
\end{document}
